\documentclass[11pt]{ctexart}
\usepackage[a4paper,margin=1in]{geometry}
\usepackage{graphicx}
\usepackage{xcolor}
\usepackage{hyperref}
\definecolor{refgray}{gray}{0.45}
\title{\textbf{市场最新观点汇总}\\[0.4em]{\large AI基础设施从GPU向CPU/内存接口芯片扩散，大宗商品定价矛盾加剧，消费与工业板块结构性分化显著}}
\author{KC桌面 自动生成}
\date{260710}
\begin{document}
\maketitle
\tableofcontents
\newpage
\section{一页摘要}
\begin{itemize}
  \item 全球宏观主线：美国经济韧性超预期但通胀回落缓慢，美联储维持双向准备；中国生产端修复与内需偏弱并存，财政传导效率成关键变量。
  \item AI基础设施：算力投资从训练向推理/Agentic AI迁移，CPU与内存接口芯片成为新受益者；BBU渗透率跃升、ABF载板转产、CoWoS扩产确认硬件升级周期未结束。
  \item 大宗商品：原油裂解价差创纪录（隐含油价超90美元/桶），霍尔木兹海峡流量恢复瓶颈在伊朗意愿；黄金定价权从央行转向ETF，铜库存降至多年低位。
  \item 行业分化：欧洲防务股从beta转向alpha，日本商社关注回购而非利润；中国电动车订单向比亚迪集中，折扣超市正结构性重塑全球零售格局。
\end{itemize}
\section{投行/券商}
AI硬件升级周期扩散至CPU/内存接口/BBU/ABF，大宗商品定价矛盾（裂解价差、霍尔木兹瓶颈）与消费分化（奢侈品定价权、折扣超市冲击）构成两大主线

\subsection{AI基础设施：从GPU单极到CPU+ASIC双轮驱动，硬件升级周期扩散}
\textbf{AI从训练转向推理和Agentic工作负载，CPU重新成为核心，带动内存接口芯片TAM 2030年达200亿美元（此前预测的3倍）；BBU渗透率从40\%跃升至85\%+，ABF载板转产加速，CoWoS扩产确认2027年放量。}

\begin{itemize}
  \item Bernstein指出，Agentic AI中CPU与GPU比例可达1:1至1:2，远高于当前训练集群的1:8，服务器CPU出货量结构性增长驱动内存接口芯片TAM CAGR从32\%跳升至65\%。
  \item Citi测算BBU在AI机架中的渗透率2025年40-45\%→2027年85\%+，单机架价值量从GB300的1.5-1.6万美元升至Rubin Ultra的3.3万美元+，电源架构正从集中式UPS向机架级分布式迁移。
  \item MS拆解2027年CoWoS产能分配，AMD MI400系列（标准版+Meta定制版）与谷歌TPU Sunfish/Zebrafish并行推进，Blackwell所谓“库存”实为供应链缓冲，2026年将被完全消化。
  \item JPM指出三星远期PE仅5.2倍，市场误判存储周期长度——NAND供应满足率仅50-60\%，企业级SSD用于AI KV缓存卸载创造全新消耗场景，结构性短缺被低估。
  \item BofA认为市场对英伟达过度悲观，当前估值已隐含2027-28年EPS 30-35\%折价，但HBM成本占比仅从5.2\%升至6.4\%，定制ASIC历史威胁被证伪。
\end{itemize}
\begin{figure}[htbp]
\centering
\includegraphics[width=0.88\linewidth]{figures/fig_001.jpg}
\caption{EXHIBIT 2 - EXHIBIT 2: As AI transitions from a training-heavy era (dominated by GPUs) to an inference-heavy era, the role of CPUs expands. Training requires parallel processing, but inference and specifically agentic workflows requires the se}
\end{figure}
\begin{figure}[htbp]
\centering
\includegraphics[width=0.88\linewidth]{figures/fig_004.jpg}
\caption{Exhibit 5 - EXHIBIT 4: The expansion of the TAM is driven by three core factors, each of which is expected to accelerate in the second half of this decade relative to the past 5 years Driver 1: Global server CPU shipment vol (Mn units)}
\end{figure}
\begin{figure}[htbp]
\centering
\includegraphics[width=0.88\linewidth]{figures/fig_761.jpg}
\caption{Figure 2 - AI Power Architecture Is Evolving, Making BBU An Increasingly Important Part of Next Gen AIDC We believe AI infrastructure is entering a new power architecture cycle as rack power density continues to increase. Historically, datacenters relied on centralized U}
\end{figure}
{\small\color{refgray}\textbf{References:} [R001] Bernstein：AI基础设施新受益者并非GPU，而是CPU与内存接口芯片; [R010] BofA：AI芯片龙头估值低于同行30\%，BofA指市场误判不确定性; [R021] JPM：三星远期市盈率仅5倍，JPM认为市场误判了存储周期长度; [R032] MS：AI芯片的“库存担忧”被高估了，2027年才是真正的放量年; [R057] Citi：不是库存回补而是架构演进，Citi详解BBU渗透率跃升路径}

\subsection{大宗商品与能源：裂解价差创纪录，霍尔木兹瓶颈在伊朗意愿，黄金定价权切换}
\textbf{原油现货与成品油价格背离达20美元（隐含油价超90美元/桶），战略储备释放人为制造价差；霍尔木兹流量恢复瓶颈不在运力而在伊朗配合意愿；黄金定价权从央行购金转向利率敏感型ETF。}

\begin{itemize}
  \item Bernstein发现裂解价差隐含油价已超90美元/桶，比现货高出逾20美元，创有记录以来最高，主因战略石油储备释放压低现货但未同步压低成品油价格。
  \item GS指出霍尔木兹海峡附近空油轮运力高达9.26亿桶，但船只不敢通行，波斯湾出口流量仅恢复至正常水平的71\%，伊朗配合意愿才是供应恢复的真正开关。
  \item JPM报告显示黄金未平仓合约周环比增长3\%，COMEX黄金净多头增加4700手至12万手，利率敏感型ETF重新获得边际定价权，预计3Q26均价4300美元/盎司。
  \item JPM指出中国铜库存过去三周累计减少8.2万吨，降至多年低位，消费改善信号明确；但原油市场陷入“中国需求变化”与“闲置产能重新入市”的双重夹击。
  \item GS下调2027年PPI预测至0\%，认为本轮能源驱动的价格冲击尾部影响将比市场预期更快消退，下半年PPI同比将逐步回落。
\end{itemize}
\begin{figure}[htbp]
\centering
\includegraphics[width=0.88\linewidth]{figures/fig_064.jpg}
\caption{Exhibit 6 - While disrupted supply chains can explain part of today's disconnect, the more important distortion is the release of government strategic reserves distorting commercial inventories (Exhibit 6) plus perhaps some ability of the Trump administration to verbally }
\end{figure}
\begin{figure}[htbp]
\centering
\includegraphics[width=0.88\linewidth]{figures/fig_067.jpg}
\caption{EXHIBIT 3 - EXHIBIT 3: Back-solving for oil price from gasoline and diesel yields good fit EXHIBIT 4: Post Covid-era...crack-implied oil price during Russia-Ukraine war was elevated \textbackslash{}\textasciitilde{}10\%...today's elevation is massive.}
\end{figure}
\begin{figure}[htbp]
\centering
\includegraphics[width=0.88\linewidth]{figures/fig_941.jpg}
\caption{Exhibit 11 - Exhibit 2: Flows Through the Strait of Hormuz Stand at \$42\textbackslash{}\%\$ of Normal Levels (7-Day Moving Average) Oil exports are estimated by taking an average of S\&P and Kpler data on the daily number of oil tankers crossing the Strait of}
\end{figure}
{\small\color{refgray}\textbf{References:} [R003] Bernstein：裂解价差创纪录，原油现货与成品油价格背离达20美元; [R012] JPM：中国铜库存降至多年低位，JPM解读消费改善信号; [R064] GS：中国6月CPI回落背后，能源和黄金价格波动才是主因; [R072] GS：特朗普撕毁停火协议后，原油供应恢复的真正瓶颈}

\subsection{消费与零售：奢侈品定价权分化，折扣超市结构性重塑，中国电动车订单集中}
\textbf{奢侈品行业“上新率”差距转化为定价权分化（香奈儿67\%领跑，Gucci被迫降价20-25\%）；折扣超市正从“低价简陋”进化为“高质低价”主流渠道，传统超市需运营重置；中国电动车7月淡季中比亚迪份额从27\%跳升至51\%。}

\begin{itemize}
  \item Bernstein手袋晴雨表显示香奈儿以67\%新品率领跑，但真正决定定价能力的是新品价格带分布——香奈儿新品集中在高端价位，而普拉达新品集中在价格带底部。
  \item BCG报告指出折扣超市在德国、挪威等成熟市场已占据超35\%份额，且正从“低价简陋”进化为“高质低价”主流渠道，高收入群体开始主动选择折扣渠道。
  \item BCG建议传统超市从顾客需求出发重置端到端流程，而非零敲碎打的降本，有望将运营利润率提升1-3个百分点。
  \item Citi/DB周度订单追踪显示比亚迪7月首周订单8.65万辆，份额从6月初27\%升至51\%，主要受海豹08上市驱动；理想/小鹏/特斯拉环比降幅均在两位数以上。
  \item JPM指出中国一线城市二手房挂牌量持续下降（较3月峰值回落3.5\%），供给侧主动调整正在为二手房价格构筑更坚实的底部。
\end{itemize}
\begin{figure}[htbp]
\centering
\includegraphics[width=0.88\linewidth]{figures/fig_071.jpg}
\caption{EXHIBIT 1 - EXHIBIT 3: ... across all regions, we now see Gucci largely straddling the pricing umbrella formed between the ‘accessible luxury’ price points (exemplified by Burberry) and the more crowded ‘core luxury’ price points occupied by P}
\end{figure}
\begin{figure}[htbp]
\centering
\includegraphics[width=0.88\linewidth]{figures/fig_073.jpg}
\caption{EXHIBIT 6 - EXHIBIT 6: Fashion media articles with static advertisements published prior to the Mercato's price change in early-May provide further evidence of Gucci's price change, with static prices listed on media outlets (left) differing f}
\end{figure}
\begin{figure}[htbp]
\centering
\includegraphics[width=0.88\linewidth]{figures/fig_1081.jpg}
\caption{EXHIBIT 1 - There is a demographic aspect to that growth: millennials prefer discounters over mainstream grocers in most markets. In large developed markets (including the US, Europe, Australia, Canada, and Japan), this population comprises 275 million people, more than a}
\end{figure}
{\small\color{refgray}\textbf{References:} [R004] Bernstein：香奈儿新品率67\%领跑，但奢侈品定价权靠的不是上新数量; [R020] JPM：从挂牌量到报价指数，中国房地产市场修复的新观察框架; [R050] Citi：不是行业回暖，而是份额重组，Citi看比亚迪如何抢回市场; [R059] DB：中国新能源车订单修复信号，比亚迪韧性凸显，蔚来修复而非反转}

\subsection{宏观与政策：美国韧性超预期但通胀双向风险，中国财政传导效率成关键}
\textbf{美国经济处于“温和但未过热”状态，美联储内部高度一致地取决于通胀回归2\%的速度；中国生产端修复与内需偏弱并存，财政存款高企说明资金尚未转化为实物工作量。}

\begin{itemize}
  \item GS解读6月FOMC纪要：几乎所有参与者同意——如果通胀“很快”消散则支持维持或降低利率，如果持续高企则支持进一步收紧，决策完全押注于通胀回归速度。
  \item GS美国宏观指标显示当前活动指标+2.6\%，二季度GDP预测维持2.2\%，金融条件收紧幅度有限，市场定价与基本面之间仍有距离。
  \item JPM指出中国5月财政存款创十年同期次高，基础设施支出同比下降12\%，专项债发行提速但资金从发行到形成实物工作量存在传导阻滞。
  \item GS认为中国央行二季度例会措辞从“稳中有进”调整为“向新向优”，首次明确提及“结构性分化”，但年内降息降准概率不大，短期政策支撑更可能来自财政支出提速。
  \item NOM模型显示人民币中间价不含逆周期因子的预测值较官方收盘价低约90个基点，逆周期因子使用已趋于克制，中间价管理正从“稳定”向“有序弹性”过渡。
  \item 另有 67 篇同来源报告纳入 references，用于校准该来源板块覆盖面。
\end{itemize}
\begin{figure}[htbp]
\centering
\includegraphics[width=0.88\linewidth]{figures/fig_903.jpg}
\caption{Exhibit 4 - Exhibit 1: Year-over-year PPI inflation edged up from May to June Exhibit 2: Food CPI inflation was roughly unchanged from May to June Exhibit 3: Goods inflation moderated mainly on lower energy prices in June}
\end{figure}
\begin{figure}[htbp]
\centering
\includegraphics[width=0.88\linewidth]{figures/fig_906.jpg}
\caption{Exhibit 3 - Exhibit 3: Goods inflation moderated mainly on lower energy prices in June Exhibit 4: PPI inflation edged up from May to June, mainly due to higher downstream prices \#\# The China Economics Team Andrew Tilton GS (Asia) L.L.C.}
\end{figure}
\begin{figure}[htbp]
\centering
\includegraphics[width=0.88\linewidth]{figures/fig_168.jpg}
\caption{Exhibit 1 - Exhibit 1: MSCI China's P/E at 10.3x is \$12\textbackslash{}\%\$ below long-term average Forward P/E valuation of MSCI China Index BofA GLOBAL RESEARCH Consensus forecast for MSCI China earnings growth BofA GLOBAL RESEARCH}
\end{figure}
{\small\color{refgray}\textbf{References:} [R009] BofA：AI主题重塑资金流向，中国消费板块观察; [R043] NOM：不是点位而是机制，NOM解读人民币中间价新信号; [R064] GS：中国6月CPI回落背后，能源和黄金价格波动才是主因; [R065] GS：央行调整例会措辞，GS解读为何短期不调整利率; [R075] GS：美联储纪要揭示双向准备，降息或加息都看通胀速度; [R079] GS：美国宏观韧性超预期，GS，资金条件收紧有限; [R086] JPM：生产端修复与需求端错位，JPM拆解中国库存周期; [R002] Bernstein：欧洲水电影响2Q26发电，唯独瑞典逆势增长; [R005] Bernstein：现金流修复才是真正拐点，Bernstein拆解防务股自证能力; [R006] Bernstein：印度信贷增长18\%韧性显现，Bernstein称宏观担忧被高估; [R007] BofA：AI蚕食创意软件护城河，Adobe的低估值是机会还是陷阱; [R008] Bernstein：功率半导体第二增长引擎，AI电力需求下的电装并购窗口; [R011] JPM：香港银行贷款增速回升，JPM提醒注意非利息收入基数效应; [R013] JPM：MLCC涨价预期落空？JPM称真正供需紧张在2027年; [R014] JPM：韩国汽车零部件公司机器人业务不再是免费期权; [R015] JPM：日本银行股上涨37\%之后，真正的分化从1Q财报开始; [R016] JPM：东南亚利润韧性强，极兔新市场跨过盈亏平衡点; [R017] 研究笔记：欧莱雅提前接手古驰美妆，利润修复预计在2028年; [R018] JPM：铝铜塑料涨价冲击汽车利润，韩元贬值对冲效果有限; [R019] JPM：JPM看辉瑞，不是价值陷阱，而是管线期权时间折价; [R022] JPM：不是复苏而是再定价，JPM拆解新加坡地产二元化行情; [R023] JPM：半导体设备利润率改善或提前兑现，JPM看好定价权驱动增长; [R024] JPM：SpaceX不只是火箭公司，JPM称其是AI基础设施新玩家; [R025] JPM：转产ABF不是临时措施，载板厂商的定价权将如何变化; [R026] JPM：特斯拉与SpaceX合并，JPM指市场低估中国变量; [R027] JPM：资源价格高位已成明牌，日本商社真正催化剂是回购计划; [R028] 关于万华化学的一份研究笔记; [R029] MS：阿里即时零售亏损收窄超预期，电商稳利润控亏损; [R030] MS：MS拆解中国医药全球化，不是卖资产，而是全球买家排队采购; [R031] JPM：散户不再一起买什么才是市场真正信号; [R033] MS：AI订单分化，高端ABF基板和MLCC成少数提价方向; [R034] MS：中国电动车7月进入淡季，但比亚迪和小鹏的局部亮点值得关注; [R035] MS：太空计算的真正起点不是数据中心，是边缘AI; [R036] MS：不是悲观而是低估，村田一季报或成消费电子观察节点; [R037] MS：人形机器人2026年中国出货5万台，供应链比整机更受益; [R038] MS：轨道边缘计算不是科幻，MS拆解供应链机会; [R039] NOM：NOM拆解日元走势，NISA散户资金流是被低估的变量; [R040] NOM：电商利润韧性超预期，NOM指阿里云正重塑公司利润结构; [R041] NOM：4亿销售额背后，中国生物制药呼吸科新品能否兑现增长？; [R042] NOM：AI不是威胁而是引擎，金蝶订阅业务规模效应显现; [R044] UBS：常规CCL涨价超预期，UBS揭示AI产能挤占下的新周期; [R045] UBS：黑石KKR利润率近70\%，另类资管如何系统蚕食市场; [R046] UBS：消费趋势变化下的免税龙头观察，中国中免收入变化需关注政策方向; [R047] BofA：市场误判Soitec光子学增速？BofA指移动端影响盈利拐点; [R048] BofA：中国医药行业交易活跃度正在驱动新一轮市场定价; [R049] Citi：面板供需平衡临界点2027年到来，公司定价权改善但新业务待验证; [R051] BofA：SpaceX垂直整合闭环，发射、应用、现金流的飞轮; [R052] Citi：CICCH股盈利加速可见性高，Citi关注机构板块估值修复; [R053] Citi：存储涨价重塑季度轮廓，AI服务器公司超预期靠的不是需求; [R054] Citi：被低估的第二曲线，Citi详解SpaceX的AI业务前景; [R055] 关于某公司近期业务动态的观察笔记; [R056] Citi：腾讯AI布局加速，混元3日均Token消耗增长20倍; [R058] DB：光伏玻璃正接近周期拐点，但多晶硅不确定性未减; [R060] DB：人形机器人收购情绪高涨，但埃斯顿利润修复可持续吗？; [R061] 从零售到批发，世纪互联利润率拐点如何兑现？; [R062] GS：GS详解一家半导体设备商，二季度超预期，毛利率先升后降再升; [R063] GS：惠普出货量下滑9\%，PC行业分化从供应链开始; [R066] GS：硬件与入境旅游跑赢软件，GS日本中小盘分化明显; [R067] 关于Costco近期运营数据的观察笔记; [R068] GS：长虹AI电视定价高420美元，GS称产品化能力成竞争关键; [R069] GS：欧莱雅接手Gucci彩妆，YSL模式能否再现; [R070] GS：Kakao盈利质量正在修复，比AI更确定的估值催化剂; [R071] GS：GS调研，强生口服新药Icotyde峰值收入或超100亿美元; [R073] GS：微软三大关注点被市场充分定价，资本开支与AI发展进入新阶段; [R074] GS：户外广告正在经历一场被低估的结构性修复; [R076] GS：西门子医疗真正问题不在估值，在诊断业务修复路径; [R077] GS：ServiceNow推100天ROI保证，AI用例从概念走向量化; [R078] GS：中国PCB外迁的真相，不是成本战，而是溢价能力; [R080] GS：GS拆解迪士尼奥兰多，客流下降2\%，消费却创历史新高; [R081] JPM：机构盈利超预期，但市场为何不买账？; [R082] 研究笔记：拜耳二季度业绩前瞻——全年指引仍是市场关注核心; [R083] 开源大模型：不是收入分流，而是能力分化的加速器; [R084] JPM：政策扩容与治理完善，人民币流动性的双重信号; [R085] JPM：中国再保险受益于海上保险回流，估值或未充分反映}

\section{战略咨询}
AI价值创造呈“门槛效应”——仅6\%企业成为领导者并创造9.3个百分点超额收益；加拿大消费韧性靠储蓄消耗和借贷支撑，折扣超市正结构性重塑全球零售格局

\subsection{AI价值创造：不是斜坡而是门槛，6\%的领导者创造全部超额收益}
\textbf{BCG研究600余家美国上市公司发现，仅6\%的企业可定义为AI应用领导者，但这6\%的企业在过去三年创造了行业调整后9.3个百分点的股东总回报超额收益，且完全由基本面（营收增长和利润率提升）贡献。}

\begin{itemize}
  \item AI领导者的价值创造路径以“增长”而非“效率”为主——59\%的领导者选择“用同样的人做更多的事”，将生产力提升重新投入业务扩张，员工人数增速比行业中位数高出3个百分点。
  \item BCG将AI能力拆解为技术栈、人才密度和部署深度三个维度，从入门到领先需经历三个关键转变：从买工具到构建专业AI生态系统、从试点到规模化部署、跨越人才鸿沟。
  \item AI领导者的收入增长和利润率提升完全由基本面驱动，而非市场情绪带来的估值扩张，说明AI价值创造是可持续的结构性改善。
\end{itemize}
\begin{figure}[htbp]
\centering
\includegraphics[width=0.88\linewidth]{figures/fig_1071.jpg}
\caption{Exhibit 1 - The methodology was developed in collaboration with Cindy Laura Stahl, a researcher at the University of St. Gallen. The scores across pillars are averaged to arrive at an overall metric. Applying our measure to a sample of publicly listed US companies, we cla}
\end{figure}
\begin{figure}[htbp]
\centering
\includegraphics[width=0.88\linewidth]{figures/fig_1072.jpg}
\caption{Exhibit 2 - The AI adoption leaders in our sample are concentrated in the tech sector, but breakout performers exist across industries. More importantly, the TSR lift cannot be dismissed as a “shovel-seller” effect: the numbers we find are industry-adjusted, meaning that }
\end{figure}
\begin{figure}[htbp]
\centering
\includegraphics[width=0.88\linewidth]{figures/fig_1073.jpg}
\caption{Exhibit 3 - AI Leaders Are Seeing a Productivity Dividend The productivity engine behind the value. Underlying this fundamental performance is a productivity advantage: revenue per employee is growing 4 percentage points faster at AI leaders than at laggards, on an indust}
\end{figure}
{\small\color{refgray}\textbf{References:} [R091] 波士顿咨询：AI价值创造不是斜坡而是门槛，波士顿咨询揭晓6\%真相}

\subsection{加拿大消费韧性：收入覆盖不足六成，储蓄消耗与借贷支撑增长}
\textbf{BCG拆解加拿大消费数据发现，2026年Q1实际消费支出增长2\%但来源正在切换——收入已不再是主要支撑，取而代之的是储蓄消耗、资产增值和借贷，中间60\%家庭正向下层靠拢。}

\begin{itemize}
  \item 中间60\%家庭的收入增长仅覆盖每新增1美元支出中的57美分，近一半新增支出来自储蓄消耗和借贷；最低20\%家庭收入增长几乎为零。
  \item 中等家庭储蓄每年减少约7000加元，低收入家庭减少15000加元；前80\%家庭资产增长13\%-26\%（股市和房产），但底层20\%资产反而缩水2\%。
  \item 中等家庭负债增长最快（+22\%），底层家庭负债增17\%但资产缩水，净资产缩水约35000加元/户；消费扩张的可持续性取决于资产价格和利率走势。
\end{itemize}
\begin{figure}[htbp]
\centering
\includegraphics[width=0.88\linewidth]{figures/fig_1076.jpg}
\caption{Exhibit 1 - \#\# Canadians Are Spending More Than They Earn Across income tiers, Canadian households raised their spending by roughly the same proportion between 2021 and 2025: up 27\% in the lowest band, 19\% in the middle, and 24\% at the top. (See Exhibit 1.) What separated}
\end{figure}
\begin{figure}[htbp]
\centering
\includegraphics[width=0.88\linewidth]{figures/fig_1077.jpg}
\caption{EXHIBIT 02 - Top 20\% \#\# EXHIBIT 02 \#\# Bottom 80\% of Earners Sliding Deeper Into Negative Savings Average yearly household savings (\textbackslash{}\$K, nominal) The middle 60\% moved from a modest positive savings position to spending more than their income, a deterioration of roughly \textbackslash{}\$7 }
\end{figure}
\begin{figure}[htbp]
\centering
\includegraphics[width=0.88\linewidth]{figures/fig_1078.jpg}
\caption{Exhibit 3 - \#\# Asset Gains Are Sustaining Spending, but Not Equally With household budgets under pressure, asset gains have enabled sustained spending across the top 80\%. These cohorts saw total assets grow by 13\% to 26\%, helped by rising market valuations. (See Exhibit 3}
\end{figure}
{\small\color{refgray}\textbf{References:} [R092] 波士顿咨询：波士顿咨询拆解加拿大消费韧性，收入覆盖不足六成}

\subsection{折扣超市崛起：从“低价简陋”到“高质低价”的主流渠道}
\textbf{BCG报告指出折扣超市正在从“宏观环境变化时的替代选择”演变为结构性的行业重塑力量，在许多地区即将占据超过一半的市场份额，传统超市的定价权出现不可逆裂缝。}

\begin{itemize}
  \item 折扣超市已打破“低价=低质”的消费者心智，高收入群体开始主动选择折扣渠道；在德国、挪威等成熟市场，折扣超市已占据超35\%份额。
  \item 传统超市需要从“卖货渠道”转变为“品牌经营者”，折扣店如Aldi、Lidl的成功证明当超市像消费品牌一样思考时，能同时实现差异化、降低成本并提高客户忠诚度。
  \item 传统超市目前的自有品牌策略是“最差的两个世界”——既没有国家品牌的资产价值，也无法在价格或品质上击败折扣店；真正的解法是从“自有品牌”升级为“零售商品牌”。
  \item 另有 2 篇同来源报告纳入 references，用于校准该来源板块覆盖面。
\end{itemize}
\begin{figure}[htbp]
\centering
\includegraphics[width=0.88\linewidth]{figures/fig_1081.jpg}
\caption{EXHIBIT 1 - There is a demographic aspect to that growth: millennials prefer discounters over mainstream grocers in most markets. In large developed markets (including the US, Europe, Australia, Canada, and Japan), this population comprises 275 million people, more than a}
\end{figure}
\begin{figure}[htbp]
\centering
\includegraphics[width=0.88\linewidth]{figures/fig_1082.jpg}
\caption{EXHIBIT 2 - \textbackslash{}- Nascent. In these markets, discounters have yet to establish a serious presence; they typically have market shares of less than 15\%. And some of these markets—including in the US, Sweden, and Australia—are still essentially white spaces. They present an opp}
\end{figure}
\begin{figure}[htbp]
\centering
\includegraphics[width=0.88\linewidth]{figures/fig_1086.jpg}
\caption{EXHIBIT 2 - Define a clear vision and strategy. The leadership team needs to establish the aspiration for how the private-label business will fit into the organization and—more important—into the overall brand architecture. Grocers can choose from three models: \textbackslash{}- The fir}
\end{figure}
{\small\color{refgray}\textbf{References:} [R093] 波士顿咨询：折扣超市占据过半市场？波士顿咨询称这不是预警而是现实; [R095] 波士顿咨询：超市行业新战场，不是价格战而是品牌战，折扣店已领先15年; [R094] 波士顿咨询：超市运营利润提升1-3个百分点的秘密，不是降本，而是重置流程; [R096] 麦肯锡：全球资产负债表显示，过去20年财富增长七成来自通胀}

\section{智库/国际机构}
IMF揭示黄金流动性远低于市场预期且避险功能退化；BIS研究显示监管成本在大银行与小银行间分布极不均匀；布鲁盖尔研究所评估欧盟关键原材料战略效果有限

\subsection{黄金储备：流动性被高估，避险功能正在退化}
\textbf{IMF最新研究报告指出，黄金是一种高不确定性储备资产，其流动性远低于名义市场价值所暗示的水平，且避险属性高度依赖于特定情境——对利率不确定性和美元贬值的对冲有效，但对市场下跌、通胀意外和地缘政治冲击的保护作用显著减弱。}

\begin{itemize}
  \item 2019-2025年央行黄金储备市值从1.2万亿飙升至4.5万亿美元，但实物量仅增长8.5\%，近三分之二增值来自金价上涨的“被动收益”而非主动买入。
  \item 黄金日均交易量约1340亿美元，但跟美债（日均1.07万亿）相比差距明显；在极端行情下的变现能力远不如账面价值。
  \item 黄金最稳定的对冲功能是针对利率不确定性和美元贬值，但2022年以来公司与债券之间传统负相关关系瓦解，进一步削弱了黄金在组合中的分散化价值。
  \item IMF建议黄金更适合放在“投资层”而非“流动性层”，央行需用流动性调整后的价值来评估实际可用资金。
\end{itemize}
\begin{figure}[htbp]
\centering
\includegraphics[width=0.88\linewidth]{figures/fig_1056.jpg}
\caption{Figure 1 - \#\# Why Are We Talking about Gold in Central Bank Reserves Now? Gold \$\textasciicircum{}\{2\}\$ has been prominent in the balance sheets of many central banks largely because of the institutional and historical legacy of the gold-standard era and subsequent postwar monetary arrang}
\end{figure}
\begin{figure}[htbp]
\centering
\includegraphics[width=0.88\linewidth]{figures/fig_1061.jpg}
\caption{Figure 3 - \#\# The Enduring Appeal: A Deep and Liquid Global Market, No Credit Risk Gold's liquidity is evidenced by substantial trading volumes in major bullion markets. Its daily average trading volume in London is about \textbackslash{}\$134 billion (Figure 3, panel 1, 2), providing c}
\end{figure}
\begin{figure}[htbp]
\centering
\includegraphics[width=0.88\linewidth]{figures/fig_1062.jpg}
\caption{Figure 3 - Figure 3. Daily Average Gold Market Trading Volume in London Panel 1: Daily volume in USD billions Panel 2: Daily volume in troy oz. millions Is gold a hedge, a diversifier, or a safe haven? Following Baur and Lucey (2010), we define a hedge as an asset that i}
\end{figure}
{\small\color{refgray}\textbf{References:} [R090] IMF：黄金不是终极避险资产？IMF揭示其流动性远低于市场预期}

\subsection{监管成本分布不均：大银行吸收，小银行被迫收紧}
\textbf{BIS工作论文研究哥伦比亚消费贷款拨备新规发现，对期限超108个月的贷款加征40\%额外拨备后，大型机构吸收了成本未转嫁借款人，而小型贷款机构明确收紧放贷。}

\begin{itemize}
  \item 大型机构在消费信贷市场占据主导份额，资产负债表足够厚，10\%或40\%的边际拨备增幅远未达到需要改变放贷策略的阈值。
  \item 小型贷款机构对108个月以上超长期贷款做出明确调整：贷款发放金额减少，LTV比例降低，说明监管成本在不同规模机构间的分布极不均匀。
  \item 窄范围、期限定向的拨备规则对整体信贷供给冲击有限，但会加剧大小银行之间的市场分化。
\end{itemize}
{\small\color{refgray}\textbf{References:} [R087] 国际清算银行：消费贷款监管加码，小银行在108个月这道门槛上退缩}

\subsection{欧盟关键原材料战略：战略项目难达自给自足目标，贸易协定比战略伙伴关系更有效}
\textbf{布鲁盖尔研究所评估欧盟关键原材料战略发现，60个战略项目产能远不足以实现2030年自给自足目标，战略伙伴关系至今未产生可测量的贸易效应，真正有效的是被低估的贸易协定条款。}

\begin{itemize}
  \item 欧盟目前全球开采份额几乎可以忽略不计，锂加工份额仅9.5\%，石墨49.4\%，距离40\%加工目标差距巨大；60个战略项目中仅21个获得明确资金支持，总金额约12.5亿欧元。
  \item 中国控制全球70\%锂精炼、90\%稀土加工，即使欧盟从别国购买矿石，最后仍需运到中国加工，这才是真正的供应链软肋。
  \item 实证分析显示，欧盟战略伙伴关系对实际贸易量提升效果不明显，反而是贸易协定中的相关条款实实在在地推动了关键原材料对欧盟的出口增长。
  \item 另有 1 篇同来源报告纳入 references，用于校准该来源板块覆盖面。
\end{itemize}
\begin{figure}[htbp]
\centering
\includegraphics[width=0.88\linewidth]{figures/fig_994.jpg}
\caption{Figure 1 - \#\# 1 Introduction As essential inputs for technologies including batteries, wind turbines, solar panels and semiconductors, critical raw materials (CRMs) are a strategic priority in European Union trade and industrial policy. Rising demand for these materials,}
\end{figure}
\begin{figure}[htbp]
\centering
\includegraphics[width=0.88\linewidth]{figures/fig_995.jpg}
\caption{Figure 2 - Figure 2: CRM demand projections in three IEA scenarios Stated policies (STEPS) -- Net Zero emissions (NZE) Announced pledges (APS) □ STEPS–NZE range a) Growth of worldwide CRM demand relative}
\end{figure}
\begin{figure}[htbp]
\centering
\includegraphics[width=0.88\linewidth]{figures/fig_996.jpg}
\caption{Figure 4 - Figure 3: Main exporters of CRMs by production stage, 2024 Figure 4: Value of EU CRM imports from extra-EU suppliers, 2024}
\end{figure}
{\small\color{refgray}\textbf{References:} [R088] 布鲁盖尔研究所：中国掌控70\%锂加工，欧盟供应链突围要靠贸易协定; [R089] IMF：财政赤字压至3\%上限，科特迪瓦如何应对中东冲击？}

\section{结语}
今日国际信源共同指向一个核心矛盾：全球宏观环境的“账面财富”在膨胀（资产价格重估），但真实的生产性投资和生产率提升并未跟上。AI硬件升级周期虽在扩散，但价值创造呈“门槛效应”——仅少数企业能跨越。大宗商品定价矛盾（裂解价差、霍尔木兹瓶颈）与消费分化（折扣超市崛起、奢侈品定价权转移）均指向结构性而非周期性的变化。
\newpage
\section{报告覆盖清单}
\begin{itemize}
  \item [R001] 投行/券商 | 投行/券商 / AI基础设施：从GPU单极到CPU+ASIC双轮驱动，硬件升级周期扩散 - Bernstein：AI基础设施新受益者并非GPU，而是CPU与内存接口芯片
  \item [R002] 投行/券商 | 投行/券商 / 宏观与政策：美国韧性超预期但通胀双向风险，中国财政传导效率成关键 - Bernstein：欧洲水电影响2Q26发电，唯独瑞典逆势增长
  \item [R003] 投行/券商 | 投行/券商 / 大宗商品与能源：裂解价差创纪录，霍尔木兹瓶颈在伊朗意愿，黄金定价权切换 - Bernstein：裂解价差创纪录，原油现货与成品油价格背离达20美元
  \item [R004] 投行/券商 | 投行/券商 / 消费与零售：奢侈品定价权分化，折扣超市结构性重塑，中国电动车订单集中 - Bernstein：香奈儿新品率67\%领跑，但奢侈品定价权靠的不是上新数量
  \item [R005] 投行/券商 | 投行/券商 / 宏观与政策：美国韧性超预期但通胀双向风险，中国财政传导效率成关键 - Bernstein：现金流修复才是真正拐点，Bernstein拆解防务股自证能力
  \item [R006] 投行/券商 | 投行/券商 / 宏观与政策：美国韧性超预期但通胀双向风险，中国财政传导效率成关键 - Bernstein：印度信贷增长18\%韧性显现，Bernstein称宏观担忧被高估
  \item [R007] 投行/券商 | 投行/券商 / 宏观与政策：美国韧性超预期但通胀双向风险，中国财政传导效率成关键 - BofA：AI蚕食创意软件护城河，Adobe的低估值是机会还是陷阱
  \item [R008] 投行/券商 | 投行/券商 / 宏观与政策：美国韧性超预期但通胀双向风险，中国财政传导效率成关键 - Bernstein：功率半导体第二增长引擎，AI电力需求下的电装并购窗口
  \item [R009] 投行/券商 | 投行/券商 / 宏观与政策：美国韧性超预期但通胀双向风险，中国财政传导效率成关键 - BofA：AI主题重塑资金流向，中国消费板块观察
  \item [R010] 投行/券商 | 投行/券商 / AI基础设施：从GPU单极到CPU+ASIC双轮驱动，硬件升级周期扩散 - BofA：AI芯片龙头估值低于同行30\%，BofA指市场误判不确定性
  \item [R011] 投行/券商 | 投行/券商 / 宏观与政策：美国韧性超预期但通胀双向风险，中国财政传导效率成关键 - JPM：香港银行贷款增速回升，JPM提醒注意非利息收入基数效应
  \item [R012] 投行/券商 | 投行/券商 / 大宗商品与能源：裂解价差创纪录，霍尔木兹瓶颈在伊朗意愿，黄金定价权切换 - JPM：中国铜库存降至多年低位，JPM解读消费改善信号
  \item [R013] 投行/券商 | 投行/券商 / 宏观与政策：美国韧性超预期但通胀双向风险，中国财政传导效率成关键 - JPM：MLCC涨价预期落空？JPM称真正供需紧张在2027年
  \item [R014] 投行/券商 | 投行/券商 / 宏观与政策：美国韧性超预期但通胀双向风险，中国财政传导效率成关键 - JPM：韩国汽车零部件公司机器人业务不再是免费期权
  \item [R015] 投行/券商 | 投行/券商 / 宏观与政策：美国韧性超预期但通胀双向风险，中国财政传导效率成关键 - JPM：日本银行股上涨37\%之后，真正的分化从1Q财报开始
  \item [R016] 投行/券商 | 投行/券商 / 宏观与政策：美国韧性超预期但通胀双向风险，中国财政传导效率成关键 - JPM：东南亚利润韧性强，极兔新市场跨过盈亏平衡点
  \item [R017] 投行/券商 | 投行/券商 / 宏观与政策：美国韧性超预期但通胀双向风险，中国财政传导效率成关键 - 研究笔记：欧莱雅提前接手古驰美妆，利润修复预计在2028年
  \item [R018] 投行/券商 | 投行/券商 / 宏观与政策：美国韧性超预期但通胀双向风险，中国财政传导效率成关键 - JPM：铝铜塑料涨价冲击汽车利润，韩元贬值对冲效果有限
  \item [R019] 投行/券商 | 投行/券商 / 宏观与政策：美国韧性超预期但通胀双向风险，中国财政传导效率成关键 - JPM：JPM看辉瑞，不是价值陷阱，而是管线期权时间折价
  \item [R020] 投行/券商 | 投行/券商 / 消费与零售：奢侈品定价权分化，折扣超市结构性重塑，中国电动车订单集中 - JPM：从挂牌量到报价指数，中国房地产市场修复的新观察框架
  \item [R021] 投行/券商 | 投行/券商 / AI基础设施：从GPU单极到CPU+ASIC双轮驱动，硬件升级周期扩散 - JPM：三星远期市盈率仅5倍，JPM认为市场误判了存储周期长度
  \item [R022] 投行/券商 | 投行/券商 / 宏观与政策：美国韧性超预期但通胀双向风险，中国财政传导效率成关键 - JPM：不是复苏而是再定价，JPM拆解新加坡地产二元化行情
  \item [R023] 投行/券商 | 投行/券商 / 宏观与政策：美国韧性超预期但通胀双向风险，中国财政传导效率成关键 - JPM：半导体设备利润率改善或提前兑现，JPM看好定价权驱动增长
  \item [R024] 投行/券商 | 投行/券商 / 宏观与政策：美国韧性超预期但通胀双向风险，中国财政传导效率成关键 - JPM：SpaceX不只是火箭公司，JPM称其是AI基础设施新玩家
  \item [R025] 投行/券商 | 投行/券商 / 宏观与政策：美国韧性超预期但通胀双向风险，中国财政传导效率成关键 - JPM：转产ABF不是临时措施，载板厂商的定价权将如何变化
  \item [R026] 投行/券商 | 投行/券商 / 宏观与政策：美国韧性超预期但通胀双向风险，中国财政传导效率成关键 - JPM：特斯拉与SpaceX合并，JPM指市场低估中国变量
  \item [R027] 投行/券商 | 投行/券商 / 宏观与政策：美国韧性超预期但通胀双向风险，中国财政传导效率成关键 - JPM：资源价格高位已成明牌，日本商社真正催化剂是回购计划
  \item [R028] 投行/券商 | 投行/券商 / 宏观与政策：美国韧性超预期但通胀双向风险，中国财政传导效率成关键 - 关于万华化学的一份研究笔记
  \item [R029] 投行/券商 | 投行/券商 / 宏观与政策：美国韧性超预期但通胀双向风险，中国财政传导效率成关键 - MS：阿里即时零售亏损收窄超预期，电商稳利润控亏损
  \item [R030] 投行/券商 | 投行/券商 / 宏观与政策：美国韧性超预期但通胀双向风险，中国财政传导效率成关键 - MS：MS拆解中国医药全球化，不是卖资产，而是全球买家排队采购
  \item [R031] 投行/券商 | 投行/券商 / 宏观与政策：美国韧性超预期但通胀双向风险，中国财政传导效率成关键 - JPM：散户不再一起买什么才是市场真正信号
  \item [R032] 投行/券商 | 投行/券商 / AI基础设施：从GPU单极到CPU+ASIC双轮驱动，硬件升级周期扩散 - MS：AI芯片的“库存担忧”被高估了，2027年才是真正的放量年
  \item [R033] 投行/券商 | 投行/券商 / 宏观与政策：美国韧性超预期但通胀双向风险，中国财政传导效率成关键 - MS：AI订单分化，高端ABF基板和MLCC成少数提价方向
  \item [R034] 投行/券商 | 投行/券商 / 宏观与政策：美国韧性超预期但通胀双向风险，中国财政传导效率成关键 - MS：中国电动车7月进入淡季，但比亚迪和小鹏的局部亮点值得关注
  \item [R035] 投行/券商 | 投行/券商 / 宏观与政策：美国韧性超预期但通胀双向风险，中国财政传导效率成关键 - MS：太空计算的真正起点不是数据中心，是边缘AI
  \item [R036] 投行/券商 | 投行/券商 / 宏观与政策：美国韧性超预期但通胀双向风险，中国财政传导效率成关键 - MS：不是悲观而是低估，村田一季报或成消费电子观察节点
  \item [R037] 投行/券商 | 投行/券商 / 宏观与政策：美国韧性超预期但通胀双向风险，中国财政传导效率成关键 - MS：人形机器人2026年中国出货5万台，供应链比整机更受益
  \item [R038] 投行/券商 | 投行/券商 / 宏观与政策：美国韧性超预期但通胀双向风险，中国财政传导效率成关键 - MS：轨道边缘计算不是科幻，MS拆解供应链机会
  \item [R039] 投行/券商 | 投行/券商 / 宏观与政策：美国韧性超预期但通胀双向风险，中国财政传导效率成关键 - NOM：NOM拆解日元走势，NISA散户资金流是被低估的变量
  \item [R040] 投行/券商 | 投行/券商 / 宏观与政策：美国韧性超预期但通胀双向风险，中国财政传导效率成关键 - NOM：电商利润韧性超预期，NOM指阿里云正重塑公司利润结构
  \item [R041] 投行/券商 | 投行/券商 / 宏观与政策：美国韧性超预期但通胀双向风险，中国财政传导效率成关键 - NOM：4亿销售额背后，中国生物制药呼吸科新品能否兑现增长？
  \item [R042] 投行/券商 | 投行/券商 / 宏观与政策：美国韧性超预期但通胀双向风险，中国财政传导效率成关键 - NOM：AI不是威胁而是引擎，金蝶订阅业务规模效应显现
  \item [R043] 投行/券商 | 投行/券商 / 宏观与政策：美国韧性超预期但通胀双向风险，中国财政传导效率成关键 - NOM：不是点位而是机制，NOM解读人民币中间价新信号
  \item [R044] 投行/券商 | 投行/券商 / 宏观与政策：美国韧性超预期但通胀双向风险，中国财政传导效率成关键 - UBS：常规CCL涨价超预期，UBS揭示AI产能挤占下的新周期
  \item [R045] 投行/券商 | 投行/券商 / 宏观与政策：美国韧性超预期但通胀双向风险，中国财政传导效率成关键 - UBS：黑石KKR利润率近70\%，另类资管如何系统蚕食市场
  \item [R046] 投行/券商 | 投行/券商 / 宏观与政策：美国韧性超预期但通胀双向风险，中国财政传导效率成关键 - UBS：消费趋势变化下的免税龙头观察，中国中免收入变化需关注政策方向
  \item [R047] 投行/券商 | 投行/券商 / 宏观与政策：美国韧性超预期但通胀双向风险，中国财政传导效率成关键 - BofA：市场误判Soitec光子学增速？BofA指移动端影响盈利拐点
  \item [R048] 投行/券商 | 投行/券商 / 宏观与政策：美国韧性超预期但通胀双向风险，中国财政传导效率成关键 - BofA：中国医药行业交易活跃度正在驱动新一轮市场定价
  \item [R049] 投行/券商 | 投行/券商 / 宏观与政策：美国韧性超预期但通胀双向风险，中国财政传导效率成关键 - Citi：面板供需平衡临界点2027年到来，公司定价权改善但新业务待验证
  \item [R050] 投行/券商 | 投行/券商 / 消费与零售：奢侈品定价权分化，折扣超市结构性重塑，中国电动车订单集中 - Citi：不是行业回暖，而是份额重组，Citi看比亚迪如何抢回市场
  \item [R051] 投行/券商 | 投行/券商 / 宏观与政策：美国韧性超预期但通胀双向风险，中国财政传导效率成关键 - BofA：SpaceX垂直整合闭环，发射、应用、现金流的飞轮
  \item [R052] 投行/券商 | 投行/券商 / 宏观与政策：美国韧性超预期但通胀双向风险，中国财政传导效率成关键 - Citi：CICCH股盈利加速可见性高，Citi关注机构板块估值修复
  \item [R053] 投行/券商 | 投行/券商 / 宏观与政策：美国韧性超预期但通胀双向风险，中国财政传导效率成关键 - Citi：存储涨价重塑季度轮廓，AI服务器公司超预期靠的不是需求
  \item [R054] 投行/券商 | 投行/券商 / 宏观与政策：美国韧性超预期但通胀双向风险，中国财政传导效率成关键 - Citi：被低估的第二曲线，Citi详解SpaceX的AI业务前景
  \item [R055] 投行/券商 | 投行/券商 / 宏观与政策：美国韧性超预期但通胀双向风险，中国财政传导效率成关键 - 关于某公司近期业务动态的观察笔记
  \item [R056] 投行/券商 | 投行/券商 / 宏观与政策：美国韧性超预期但通胀双向风险，中国财政传导效率成关键 - Citi：腾讯AI布局加速，混元3日均Token消耗增长20倍
  \item [R057] 投行/券商 | 投行/券商 / AI基础设施：从GPU单极到CPU+ASIC双轮驱动，硬件升级周期扩散 - Citi：不是库存回补而是架构演进，Citi详解BBU渗透率跃升路径
  \item [R058] 投行/券商 | 投行/券商 / 宏观与政策：美国韧性超预期但通胀双向风险，中国财政传导效率成关键 - DB：光伏玻璃正接近周期拐点，但多晶硅不确定性未减
  \item [R059] 投行/券商 | 投行/券商 / 消费与零售：奢侈品定价权分化，折扣超市结构性重塑，中国电动车订单集中 - DB：中国新能源车订单修复信号，比亚迪韧性凸显，蔚来修复而非反转
  \item [R060] 投行/券商 | 投行/券商 / 宏观与政策：美国韧性超预期但通胀双向风险，中国财政传导效率成关键 - DB：人形机器人收购情绪高涨，但埃斯顿利润修复可持续吗？
  \item [R061] 投行/券商 | 投行/券商 / 宏观与政策：美国韧性超预期但通胀双向风险，中国财政传导效率成关键 - 从零售到批发，世纪互联利润率拐点如何兑现？
  \item [R062] 投行/券商 | 投行/券商 / 宏观与政策：美国韧性超预期但通胀双向风险，中国财政传导效率成关键 - GS：GS详解一家半导体设备商，二季度超预期，毛利率先升后降再升
  \item [R063] 投行/券商 | 投行/券商 / 宏观与政策：美国韧性超预期但通胀双向风险，中国财政传导效率成关键 - GS：惠普出货量下滑9\%，PC行业分化从供应链开始
  \item [R064] 投行/券商 | 投行/券商 / 大宗商品与能源：裂解价差创纪录，霍尔木兹瓶颈在伊朗意愿，黄金定价权切换 / 投行/券商 / 宏观与政策：美国韧性超预期但通胀双向风险，中国财政传导效率成关键 - GS：中国6月CPI回落背后，能源和黄金价格波动才是主因
  \item [R065] 投行/券商 | 投行/券商 / 宏观与政策：美国韧性超预期但通胀双向风险，中国财政传导效率成关键 - GS：央行调整例会措辞，GS解读为何短期不调整利率
  \item [R066] 投行/券商 | 投行/券商 / 宏观与政策：美国韧性超预期但通胀双向风险，中国财政传导效率成关键 - GS：硬件与入境旅游跑赢软件，GS日本中小盘分化明显
  \item [R067] 投行/券商 | 投行/券商 / 宏观与政策：美国韧性超预期但通胀双向风险，中国财政传导效率成关键 - 关于Costco近期运营数据的观察笔记
  \item [R068] 投行/券商 | 投行/券商 / 宏观与政策：美国韧性超预期但通胀双向风险，中国财政传导效率成关键 - GS：长虹AI电视定价高420美元，GS称产品化能力成竞争关键
  \item [R069] 投行/券商 | 投行/券商 / 宏观与政策：美国韧性超预期但通胀双向风险，中国财政传导效率成关键 - GS：欧莱雅接手Gucci彩妆，YSL模式能否再现
  \item [R070] 投行/券商 | 投行/券商 / 宏观与政策：美国韧性超预期但通胀双向风险，中国财政传导效率成关键 - GS：Kakao盈利质量正在修复，比AI更确定的估值催化剂
  \item [R071] 投行/券商 | 投行/券商 / 宏观与政策：美国韧性超预期但通胀双向风险，中国财政传导效率成关键 - GS：GS调研，强生口服新药Icotyde峰值收入或超100亿美元
  \item [R072] 投行/券商 | 投行/券商 / 大宗商品与能源：裂解价差创纪录，霍尔木兹瓶颈在伊朗意愿，黄金定价权切换 - GS：特朗普撕毁停火协议后，原油供应恢复的真正瓶颈
  \item [R073] 投行/券商 | 投行/券商 / 宏观与政策：美国韧性超预期但通胀双向风险，中国财政传导效率成关键 - GS：微软三大关注点被市场充分定价，资本开支与AI发展进入新阶段
  \item [R074] 投行/券商 | 投行/券商 / 宏观与政策：美国韧性超预期但通胀双向风险，中国财政传导效率成关键 - GS：户外广告正在经历一场被低估的结构性修复
  \item [R075] 投行/券商 | 投行/券商 / 宏观与政策：美国韧性超预期但通胀双向风险，中国财政传导效率成关键 - GS：美联储纪要揭示双向准备，降息或加息都看通胀速度
  \item [R076] 投行/券商 | 投行/券商 / 宏观与政策：美国韧性超预期但通胀双向风险，中国财政传导效率成关键 - GS：西门子医疗真正问题不在估值，在诊断业务修复路径
  \item [R077] 投行/券商 | 投行/券商 / 宏观与政策：美国韧性超预期但通胀双向风险，中国财政传导效率成关键 - GS：ServiceNow推100天ROI保证，AI用例从概念走向量化
  \item [R078] 投行/券商 | 投行/券商 / 宏观与政策：美国韧性超预期但通胀双向风险，中国财政传导效率成关键 - GS：中国PCB外迁的真相，不是成本战，而是溢价能力
  \item [R079] 投行/券商 | 投行/券商 / 宏观与政策：美国韧性超预期但通胀双向风险，中国财政传导效率成关键 - GS：美国宏观韧性超预期，GS，资金条件收紧有限
  \item [R080] 投行/券商 | 投行/券商 / 宏观与政策：美国韧性超预期但通胀双向风险，中国财政传导效率成关键 - GS：GS拆解迪士尼奥兰多，客流下降2\%，消费却创历史新高
  \item [R081] 投行/券商 | 投行/券商 / 宏观与政策：美国韧性超预期但通胀双向风险，中国财政传导效率成关键 - JPM：机构盈利超预期，但市场为何不买账？
  \item [R082] 投行/券商 | 投行/券商 / 宏观与政策：美国韧性超预期但通胀双向风险，中国财政传导效率成关键 - 研究笔记：拜耳二季度业绩前瞻——全年指引仍是市场关注核心
  \item [R083] 投行/券商 | 投行/券商 / 宏观与政策：美国韧性超预期但通胀双向风险，中国财政传导效率成关键 - 开源大模型：不是收入分流，而是能力分化的加速器
  \item [R084] 投行/券商 | 投行/券商 / 宏观与政策：美国韧性超预期但通胀双向风险，中国财政传导效率成关键 - JPM：政策扩容与治理完善，人民币流动性的双重信号
  \item [R085] 投行/券商 | 投行/券商 / 宏观与政策：美国韧性超预期但通胀双向风险，中国财政传导效率成关键 - JPM：中国再保险受益于海上保险回流，估值或未充分反映
  \item [R086] 投行/券商 | 投行/券商 / 宏观与政策：美国韧性超预期但通胀双向风险，中国财政传导效率成关键 - JPM：生产端修复与需求端错位，JPM拆解中国库存周期
  \item [R087] 智库/国际机构 | 智库/国际机构 / 监管成本分布不均：大银行吸收，小银行被迫收紧 - 国际清算银行：消费贷款监管加码，小银行在108个月这道门槛上退缩
  \item [R088] 智库/国际机构 | 智库/国际机构 / 欧盟关键原材料战略：战略项目难达自给自足目标，贸易协定比战略伙伴关系更有效 - 布鲁盖尔研究所：中国掌控70\%锂加工，欧盟供应链突围要靠贸易协定
  \item [R089] 智库/国际机构 | 智库/国际机构 / 欧盟关键原材料战略：战略项目难达自给自足目标，贸易协定比战略伙伴关系更有效 - IMF：财政赤字压至3\%上限，科特迪瓦如何应对中东冲击？
  \item [R090] 智库/国际机构 | 智库/国际机构 / 黄金储备：流动性被高估，避险功能正在退化 - IMF：黄金不是终极避险资产？IMF揭示其流动性远低于市场预期
  \item [R091] 战略咨询 | 战略咨询 / AI价值创造：不是斜坡而是门槛，6\%的领导者创造全部超额收益 - 波士顿咨询：AI价值创造不是斜坡而是门槛，波士顿咨询揭晓6\%真相
  \item [R092] 战略咨询 | 战略咨询 / 加拿大消费韧性：收入覆盖不足六成，储蓄消耗与借贷支撑增长 - 波士顿咨询：波士顿咨询拆解加拿大消费韧性，收入覆盖不足六成
  \item [R093] 战略咨询 | 战略咨询 / 折扣超市崛起：从“低价简陋”到“高质低价”的主流渠道 - 波士顿咨询：折扣超市占据过半市场？波士顿咨询称这不是预警而是现实
  \item [R094] 战略咨询 | 战略咨询 / 折扣超市崛起：从“低价简陋”到“高质低价”的主流渠道 - 波士顿咨询：超市运营利润提升1-3个百分点的秘密，不是降本，而是重置流程
  \item [R095] 战略咨询 | 战略咨询 / 折扣超市崛起：从“低价简陋”到“高质低价”的主流渠道 - 波士顿咨询：超市行业新战场，不是价格战而是品牌战，折扣店已领先15年
  \item [R096] 战略咨询 | 战略咨询 / 折扣超市崛起：从“低价简陋”到“高质低价”的主流渠道 - 麦肯锡：全球资产负债表显示，过去20年财富增长七成来自通胀
\end{itemize}
\section{Disclaimer}
{\scriptsize\color{refgray}
This community is a paid learning and information-sharing community created for general educational purposes only. The membership fee is charged solely for access to community discussions, educational content, organization, and operational costs. It is not an advisory fee, management fee, performance fee, brokerage fee, or compensation for personalized financial, investment, legal, tax, or professional advice.\par
I participate and share information solely as an individual and for educational discussion among community members. I am not acting as a financial adviser, investment adviser, broker, dealer, portfolio manager, legal adviser, tax adviser, accountant, fiduciary, or any other licensed professional adviser. Nothing shared in this community constitutes, and should not be interpreted as, financial advice, investment advice, legal advice, tax advice, a recommendation, solicitation, offer, endorsement, instruction, or guarantee to buy, sell, hold, short, trade, or invest in any security, cryptocurrency, fund, derivative, strategy, or other financial product.\par
All content shared in this community, including messages, discussions, links, files, charts, examples, market commentary, watchlists, AI-generated or AI-polished summaries, and third-party materials, is general, impersonal, and educational in nature. It does not take into account any individual's financial situation, investment objectives, risk tolerance, investment horizon, liquidity needs, tax status, legal restrictions, or suitability requirements. No content should be relied upon as suitable for any specific person, account, portfolio, or financial decision.\par
Information shared in this community may be incomplete, inaccurate, outdated, speculative, AI-generated, AI-edited, or based on third-party internet sources that have not been independently verified. I make no representation or warranty as to the accuracy, completeness, timeliness, reliability, or suitability of any information shared.\par
Financial markets involve substantial risk, including the possible loss of principal. Past performance, historical data, hypothetical examples, simulated results, backtests, personal opinions, or educational examples do not guarantee future results. Each member is solely responsible for conducting their own research, exercising independent judgment, and making their own decisions. Before making any financial, investment, legal, or tax decision, members should consult qualified and licensed professionals.\par
Participation in this community, payment of any membership fee, or communication with me or other members does not create any adviser-client, fiduciary, professional, contractual, agency, partnership, employment, or other advisory relationship. By joining or participating in this community, you acknowledge and agree that you are solely responsible for your own decisions, actions, transactions, profits, losses, risks, and consequences, and that I shall not be liable for any direct, indirect, incidental, consequential, financial, legal, tax, or other loss or damage arising from or related to any content, discussion, information, or material shared in this community.\par
}
\end{document}
